\section{总结与延伸}

\subsection{核心要点回顾}

本讲从三个角度深入阐述了推理时引导生成的理论与实践：

\begin{enumerate}
    \item \textbf{SDE/ODE 视角的统一}：DPS 引导不仅适用于 DDPM 的离散表述，也可以在 SDE 和 ODE 的连续时间表述中一致应用。
    \item \textbf{RL 视角的重新诠释}：通过 KL 约束的奖励最大化框架，推导出了与 DPS 等价的引导公式。这一 RL 视角揭示了推理时引导与 RLHF 的深层联系。
    \item \textbf{丰富的应用场景}：通过 canonical space/instance space 的同步化技术，预训练模型可以生成全景图、3D 纹理、歧义图像等远超训练域的内容。
\end{enumerate}

\subsection{与 LLM Test-Time Scaling 的呼应}

本讲将扩散模型的推理时引导置于更宏大的 AI 趋势中——test-time scaling。LLM 领域通过 O1/O3/DeepSeek R1 展示了在推理时花费更多计算的巨大潜力。扩散模型的推理时引导本质上也是一种 test-time scaling 技术，它不增加训练成本，而是通过推理时的额外计算来提升输出质量。

\subsection{总结表：方法、价值与风险}

\begin{table}[H]
\centering
\begin{tabular}{p{0.22\textwidth} p{0.21\textwidth} p{0.23\textwidth} p{0.22\textwidth}}
\toprule
主题 & 核心价值 & 主要风险 & 实践建议 \\
\midrule
DPS / 逆问题引导 & 无需重训即可处理观测约束 & 高噪声阶段近似误差大 & 采用较多步数并做 guidance schedule \\
RL 视角奖励引导 & 把问题统一到 verifier + reward 框架 & reward hacking、价值近似粗糙 & 先做 verifier 校验，再谈更强引导 \\
同步化生成 & 用 2D 预训练模型外推到 panorama / 3D / illusion & 多实例间约束冲突，易失真 & 在 canonical / instance space 间显式记录投影与重叠区域 \\
Test-time scaling & 以较低训练成本换更高推理质量 & 计算开销上升，延迟更高 & 用 best-of-$N$ 或自适应算力分配控制成本 \\
\bottomrule
\end{tabular}
\caption{Lecture 12 的方法总览}
\end{table}

\subsection{给研究者的三个行动点}

\begin{enumerate}
    \item 先在一个简单、可微、可评测的 reward 上建立稳定基线，再扩展到多约束设置
    \item 为每种方法建立统一的失败案例库，特别记录高引导权重与低步数下的退化模式
    \item 把扩散模型的推理时引导和 LLM 的 verifier / reranking / adaptive compute 方法对照研究，而不是分开看待
\end{enumerate}

\subsection{进一步延伸的问题}

\begin{itemize}
    \item 是否能训练一个统一 verifier，同时服务逆问题、编辑任务与美学偏好？
    \item 是否可以让模型动态决定哪些样本值得更多推理步数，从而做到真正的 adaptive test-time scaling？
    \item 对于 3D、视频和机器人轨迹等更复杂模态，canonical space / instance space 的同步化是否还能保持稳定？
\end{itemize}

\subsection{与训练时控制方法的分工}

把本讲内容放回更完整的生成系统视角，可以得到一个很实用的判断：推理时引导并不是要取代 ControlNet、LoRA 或条件扩散模型，而是与它们分工协作。

\begin{table}[H]
\centering
\begin{tabular}{p{0.22\textwidth} p{0.24\textwidth} p{0.22\textwidth} p{0.2\textwidth}}
\toprule
路线 & 最擅长的问题 & 主要代价 & 更适合的场景 \\
\midrule
训练时控制 & 固定且高频的条件类型 & 需要数据、训练和维护多个模型 & 稳定生产环境、低延迟服务 \\
推理时引导 & 临时、多变、组合式约束 & 采样更慢、调参更敏感 & 研究探索、长尾需求、复杂约束拼装 \\
混合路线 & 训练时给粗结构，推理时做精修 & 系统设计更复杂 & 高质量生成与灵活控制同时要求较高的任务 \\
\bottomrule
\end{tabular}
\caption{训练时控制与推理时引导的职责边界}
\end{table}

这张表的核心启示是：如果某种条件几乎每次都出现，最好把它吸收到训练阶段；如果条件是临时的、跨任务的、难以收集监督数据的，就更适合保留到推理阶段解决。

\subsection{拓展阅读}

\begin{itemize}
    \item \textbf{DPS 原始论文}：Chung et al., ``Diffusion Posterior Sampling for General Noisy Inverse Problems,'' ICLR 2023
    \item \textbf{Test-time scaling survey}：讲者推荐了一篇关于 test-time scaling 技术的综述论文，涵盖了扩散模型和流模型领域的多种技术
    \item \textbf{SyncDiffusion}：Lee et al., ``SyncDiffusion: Coherent Montage via Synchronized Joint Diffusions,'' NeurIPS 2023——全景图像同步化生成
    \item \textbf{Visual Anagrams}：Geng et al., ``Visual Anagrams: Generating Multi-View Optical Illusions with Diffusion Models,'' CVPR 2024——歧义图像生成
    \item \textbf{TEXTure}：Richardson et al., ``TEXTure: Text-Guided Texturing of 3D Meshes,'' SIGGRAPH 2023——利用 2D 扩散模型进行 3D 纹理生成
    \item \textbf{RLHF for Diffusion}：Fan et al., ``DPOK: Reinforcement Learning for Fine-tuning Text-to-Image Diffusion Models''——将 RLHF 思想引入扩散模型训练
\end{itemize}

%% --- Fin del contenido --- %%

\end{document}
